\begin{afterword}
\summaryhead

The essay asks what separates a technical understanding from a verbal one, and answers with probability theory. A belief is a distribution of probability over future outcomes, a fixed amount that must be shared out. To make forecasters report their real probabilities, bets must be scored by a proper rule. Requiring that the total score not depend on how experiments are grouped leaves only one rule: the logarithm of the probability given to what happened. That score rewards calibration, saying 90 per cent when you are right nine times in ten, and also discrimination, telling true answers from false.

Worked examples follow. A precise hypothesis that bets 90 per cent on the result 51 gains a tenfold advantage over a vague one when 51 appears; a hypothesis that fits everything is no better than ignorance. Human reasoning, by contrast, runs on a feeling of ``fit'' that is not conserved and is judged after the fact. That is the difference between verbal and technical explanation.

The essay then qualifies its standard. Nineteenth-century evolution was vague, yet right against precise physics about the age of the Sun. Advance prediction is a social rule of science, not a law of probability. Critics of a confirmed theory must say exactly what it predicts badly. Popper's falsifiability becomes a quantitative matter. A parable of ``Bayesianity'' shows why no probability should be 1, and a waking blue tentacle shows that one cannot explain an event one does not in advance anticipate.

\responsehead

The mathematical core of the essay is correct. I checked every formula and worked number by script, and the working is in the report. The logarithm does follow from requiring that scores add when probabilities multiply. The precise theory's tenfold advantage, the fixed coin, the three-result scores in decibels and the sunrise figures are all correct, apart from roundings that change nothing. Some of the ideas are older than the essay and are not credited to their authors: the logarithmic rule is I. J. Good's (1952), and the split of a score into calibration and discrimination is Murphy's (1973).

The one technical error is in the scoring rules. The essay offers ``a dollar minus the squared error of the bet'' on the winning colour as a proper rule, in a game with three colours. It is not proper there. Facing lights at 30, 20 and 50 per cent, a bettor earns more on average by betting about 35, 3 and 61 cents than by betting the true frequencies. The footnote that checks the rule covers two colours only. Brier's rule, which charges the error on every colour, is proper, and once payment depends only on the winning bet, the logarithm is the only strictly proper rule for three or more outcomes. That supports the essay's choice of the logarithm, by a different route from the one it gives.

The rival positions are described more narrowly than their own sources describe them. Classical statistics is said to attend only to likelihoods; the standard complaint against it is that it violates the likelihood principle, and it has its own penalties for fitted hypotheses. Popper's call for bold theories is translated into likelihoods, although Popper preferred the more improbable theory, as the notes on ``Science Isn't Strict Enough'' explain. Defending a hypothesis is said to have ``no analogue'' in probability or decision theory; Lakatos's account of research programmes treats such defence as rational, and its standard example is the search for an unseen planet that the essay itself tells as Newton's triumph. On advance prediction the essay takes the side of Mill and Keynes against Whewell and Maher without saying there is a dispute, and its claim that explanation after the fact ``was correct as probability theory'' faces the problem of old evidence, which it does not mention.

Several historical claims are overstated, though the arguments still hold after correction. Halley's comet confirmed an advance prediction of Newton's theory in 1758, long before Neptune. The first edition of the \textsc{Origin} puts three hundred million years on the erosion of the Weald, so nineteenth-century evolution had at least a scrap of quantitative argument. The figure of one part in $10^{14}$ for relativity could be traced only to Penrose; the best-known pulsar test agrees to about 0.2 per cent.

The section on the blue tentacle overstates its point. Explaining an event one does not expect is said to ``necessarily'' violate probability theory. But asking what would best explain an unlikely event is an ordinary conditional question, and the essay's own count of two million beings answers it. What probability forbids is holding an explanation probable without expecting what it predicts.

In short: a correct account of logarithmic scoring and of why precise hypotheses win, with one rule called proper that is not, and with the positions it argues against (classical statistics, Popper on probability, Lakatos on defending theories, predictivism) described more narrowly than their sources describe them.
\end{afterword}
